
There are two ways to use the library. The first one is to use it as
a new graph data-type for your graph algorithms. The second one is to
bind data to a graph. The second way is normally used if you have
any data which skeleton structure is a graph, e.g. data-flow
graphs. The library helps you to maintain and work with your data
and gives you powerful features to display the data.  The main
classes to access the library are {\bf FFGraph}, {\bf
  FFGraphClass}, {\bf FFNodeClass} and {\bf FFEdgeClass}. In
addition there are two data types {\bf FFNode} and {\bf FFEdge}
which are pointers to {\bf FFNodeClass} and {\bf FFEdgeClass}.
{\bf FFGraph} is a container for {\bf FFGraphClass} objects. The
user has no direct access to the class {\bf FFGraphClass}. The
user can call methods on {\bf FFGraphClass} objects by using the
\verb|->| operator on {\bf FFGraph} objects. E.g. to create a new
graph and add a node to the graph you would write:

\bigskip

{\tt  
FFGraph g; \\
g\verb|->|addNode(); \\    
}

\bigskip

The reason for this is the overloaded \verb|->| operator in {\bf
  FFGraph} which passes all method calls to the underlying {\bf
  FFGraphClass} object. Copies of {\bf FFGraph} objects share the
same {\bf FFGraphClass} object. This allows an efficient garbage
collection and protects the user from making any of the usual pointer
errors like uninitialized pointers or pointers to invalid addresses. If
you want a copy of a graph you can use the method {\bf
FFGraphClass::copy(...)}. However {\bf FFNode} and {\bf FFEdge} are
only pointers, for efficiency reasons so beware\footnote{Using
  the container mechanism for FFNode and FFEdge would slow down
  the library by 30\%}!  For a detailed description of all
methods of the classes {\bf FFGraph}, {\bf FFGraphClass}, {\bf
  FFNodeClass} and {\bf FFEdgeClass} see
sections\vref{standard_classes}.

The following data types which are not C++ classes are defined in
the library. 

\begin{itemize}

\item{\bf BOOL}
  
  {\bf BOOL} is a boolean data type with the two values {\bf YES}
  and {\bf NO}. 

\item{\bf FFCoord}

  {\bf FFCoord} is a floating point data type.
  
\end{itemize}

If you have a label that should be bound to every graph, node or 
edge, you can use the methods

\begin{itemize}

\item{\bf void addDefGraphLabel(FFGraphLabel *)}

\item{\bf void addDefNodeLabel(FFNodeLabel *)}

\item{\bf void addDefEdgeLabel(FFEdgeLabel *)}

\end{itemize}

to specify default labels for them.\\

\bigskip

The library uses two kinds of error handlers:

\begin{itemize}

\item{\bf ffError(errGrade, errKind, char *file, int line)}

\item{\bf ffError(errGrade, char *msg, char *file, int line)}

\end{itemize}

 The library calls these functions whenever an error {\bf kind} or
{\bf msg} of grade {\bf grade} occurred in file {\bf file} at line
{\bf line}. The values of {\bf grade} and {\bf kind} are defined in
{\bf ffError.h}.  If you want to display errors differently, you can
tell the library to use your own error-handler by calling:

\begin{itemize}

\item{\bf setErrorHandler(kindErrorHandler, charErrorHandler)}

{\bf kindErrorHandler} and {\bf charErrorHandler} are defined in {\bf
ffError.h}. If you pass {\bf NULL} as a handler, the handler remains
unchanged.

\end{itemize}
